\sectionheading{5.27 Programming Environment}

কোনো Programming Language-এ Program লিখে সেটি চালানোর জন্য
শুধু ভাষাটি জানলেই হয় না।

Code লেখা, Edit করা, Run করা এবং Error শনাক্ত ও সংশোধন করার জন্য
প্রয়োজনীয় Software ও Tool-এর সমন্বিত পরিবেশ প্রয়োজন হয়।
এই পরিবেশকে \textbf{Programming Environment} বলা হয়।

%================================================
% Definition
%================================================

\booksubheading{Programming Environment কাকে বলে?}

\begin{definitionbox}

\textbf{\color{darkTeal}সংজ্ঞা}

Program \textbf{লেখা, সম্পাদনা (Edit), Execute এবং Error
শনাক্ত ও সংশোধন} করার জন্য প্রয়োজনীয় Software, Tool ও
সুবিধার সমন্বিত ব্যবস্থাকে \textbf{Programming Environment} বলে।

\end{definitionbox}

\textbf{সহজভাবে মনে রাখবে—}

\begin{center}

\textbf{Code লেখা}
\(\longrightarrow\)
\textbf{Run করা}
\(\longrightarrow\)
\textbf{Error দেখা}
\(\longrightarrow\)
\textbf{সংশোধন করা}

\end{center}

এই কাজগুলো করার প্রয়োজনীয় পরিবেশই হলো
\textbf{Programming Environment}।

%================================================
% Main Components
%================================================

\booksubheading{Python Programming Environment-এর প্রধান উপাদান}

Python Program তৈরি ও চালানোর জন্য সাধারণত নিচের
উপাদানগুলো ব্যবহৃত হয়—

\begin{tabularx}
{\linewidth}
{
|>{\bfseries\RaggedRight\arraybackslash}p{0.32\linewidth}
|Y|
}

\hline

\rowcolor{softBlue}

\multicolumn{1}{|c|}{\textbf{উপাদান}}
&
\multicolumn{1}{c|}{\textbf{কাজ}}

\\

\hline

Python Interpreter
&
Python Code বুঝে Execute করতে সাহায্য করে।
\\

\hline

Code Editor
&
Program Code লেখা, Edit ও Save করা হয়।
\\

\hline

Shell / Console
&
সরাসরি Instruction দিয়ে তাৎক্ষণিক Result দেখা যায়।
\\

\hline

IDE
&
Code লেখা, Run, Debug এবং অন্যান্য Development সুবিধা
এক জায়গায় পাওয়া যায়।
\\

\hline

\end{tabularx}

% ---------- Figure: Programming Environment ----------
\replaceableimage[60mm]
{images/python-environment.png}
{Programming Environment-এর Code Editor, Run ও Output/Terminal অংশ}

\begin{center}
\footnotesize\textbf{চিত্র: একটি Programming Environment-এর প্রধান অংশ}
\end{center}

%================================================
% Interpreter
%================================================

\booksubheading{১. Python Interpreter}

Python Interpreter হলো এমন Software যা Python Program-এর
Instruction বুঝে Execute করতে সাহায্য করে।

উদাহরণ—

\begin{codebox}

\begin{english}
\ttfamily
print("Hello")
\end{english}

\end{codebox}

এই Python Code Execute করার জন্য Python Interpreter প্রয়োজন।

\begin{center}

\textbf{Python Code}
\(\longrightarrow\)
\textbf{Interpreter}
\(\longrightarrow\)
\textbf{Execution}
\(\longrightarrow\)
\textbf{Output}

\end{center}

\begin{rememberbox}

\textbf{\color{deepOrange}MCQ তথ্য}

\textbf{Interpreter}
\(\longrightarrow\)
Python Code Execute করতে সাহায্য করে।

\end{rememberbox}

%================================================
% Code Editor
%================================================

\booksubheading{২. Code Editor}

যে Software-এ Program Code লেখা ও পরিবর্তন করা যায় তাকে
\textbf{Code Editor} বলে।

Code Editor ব্যবহার করে—

\begin{itemize}

\item নতুন Program Code লেখা যায়।

\item Code Edit বা পরিবর্তন করা যায়।

\item Program File Save করা যায়।

\item Code সুন্দর ও সুশৃঙ্খলভাবে লেখা যায়।

\end{itemize}

উদাহরণ হিসেবে
\textbf{Visual Studio Code (VS Code)}
-এর মতো Code Editor ব্যবহার করা যায়।

\begin{rememberbox}

Code Editor-এর বিভিন্ন মাধ্যম ও ব্যবহার
\textbf{5.29 Python-এ Code লেখার মাধ্যম}
অংশে বিস্তারিত আলোচনা করা হবে।

\end{rememberbox}

%================================================
% Shell
%================================================

\booksubheading{৩. Python Shell / Interactive Mode}

Python Shell-এ সাধারণত একটি Instruction লিখে Enter চাপলে
সঙ্গে সঙ্গে তার Result দেখা যায়।

\begin{examplebox}{উদাহরণ}

\begin{codebox}

\begin{english}
\ttfamily
>>> 10 + 20\\
30
\end{english}

\end{codebox}

এখানে
\textenglish{\ttfamily 10 + 20}
Instruction দেওয়ার পর সঙ্গে সঙ্গে
\textenglish{\ttfamily 30}
Result পাওয়া গেছে।

\end{examplebox}

\begin{center}

\textbf{একটি Instruction}
\(\longrightarrow\)
\textbf{সঙ্গে সঙ্গে Result}

\end{center}

ছোট Calculation বা কোনো ছোট Code দ্রুত পরীক্ষা করার জন্য
Shell সুবিধাজনক।

%================================================
% IDE
%================================================

\booksubheading{৪. IDE — Integrated Development Environment}

\begin{definitionbox}

\textbf{\color{darkTeal}IDE-এর পূর্ণরূপ}

\textbf{Integrated Development Environment}

\end{definitionbox}

IDE এমন একটি Software Environment যেখানে Program Development-এর
বিভিন্ন সুবিধা একই জায়গায় পাওয়া যায়।

IDE-তে সাধারণত—

\begin{itemize}

\item Code Editor থাকে।

\item Program Run/Execute করার সুবিধা থাকে।

\item Error Message দেখা যায়।

\item Debugging-এর সুবিধা থাকে।

\item Program File ব্যবস্থাপনার সুবিধা থাকে।

\end{itemize}

\textbf{সহজভাবে মনে রাখবে—}

\begin{center}

\textbf{Editor}
+
\textbf{Run}
+
\textbf{Debug}
+
\textbf{অন্যান্য Tool}
=
\textbf{IDE}

\end{center}

%================================================
% Common Environments
%================================================

\booksubheading{Python-এর পরিচিত Programming Environment ও Tool}

Python Code লেখা ও চালানোর জন্য বিভিন্ন Environment বা Tool
ব্যবহার করা যায়। যেমন—

\begin{itemize}

\item \textbf{IDLE}

\item \textbf{Visual Studio Code (VS Code)}

\item \textbf{PyCharm}

\item \textbf{Jupyter Notebook}

\end{itemize}

এগুলোর কাজের ধরন ও সুবিধা এক নয়, তবে Python Program
লেখা ও চালানোর কাজে ব্যবহার করা যায়।

%================================================
% IDLE
%================================================

\booksubheading{IDLE}

Python শেখার ক্ষেত্রে পরিচিত একটি Environment হলো
\textbf{IDLE}।

\begin{definitionbox}

\textbf{\color{darkTeal}IDLE-এর পূর্ণরূপ}

\textbf{Integrated Development and Learning Environment}

\end{definitionbox}

IDLE-এর মাধ্যমে—

\begin{itemize}

\item Python Code লেখা যায়।

\item Python Shell ব্যবহার করা যায়।

\item Program Run করা যায়।

\item Program File Save করা যায়।

\item Error Message দেখা যায়।

\end{itemize}

\begin{rememberbox}

\textbf{\color{deepOrange}MCQ তথ্য}

\textbf{IDE}
\(\rightarrow\)
Integrated Development Environment

\medskip

\textbf{IDLE}
\(\rightarrow\)
Integrated Development and Learning Environment

\end{rememberbox}

%================================================
% Shell vs Editor
%================================================

\booksubheading{Shell ও Editor-এর পার্থক্য}

\begin{tabularx}
{\linewidth}
{
|>{\bfseries\RaggedRight\arraybackslash}p{0.20\linewidth}
|Y
|Y|
}

\hline

\rowcolor{softBlue}

\multicolumn{1}{|c|}{\textbf{বিষয়}}
&
\multicolumn{1}{c|}{\textbf{Shell}}
&
\multicolumn{1}{c|}{\textbf{Editor}}

\\

\hline

কাজ
&
সরাসরি Instruction Execute করা হয়।
&
পূর্ণ Program লেখা ও Edit করা হয়।
\\

\hline

Result
&
সাধারণত সঙ্গে সঙ্গে Result দেখা যায়।
&
Program Run করার পর Result দেখা যায়।
\\

\hline

ব্যবহার
&
ছোট Test বা Calculation-এর জন্য সুবিধাজনক।
&
বড় বা সংরক্ষণযোগ্য Program লেখার জন্য সুবিধাজনক।
\\

\hline

File
&
প্রতিবার আলাদা File তৈরি করা প্রয়োজন হয় না।
&
Program Code File হিসেবে Save করা যায়।
\\

\hline

\end{tabularx}

\begin{rememberbox}

\textbf{\color{deepOrange}সহজভাবে মনে রাখবে}

\textbf{Shell}
\(\longrightarrow\)
Quick Test / Quick Result

\medskip

\textbf{Editor}
\(\longrightarrow\)
Write, Edit and Save Program

\end{rememberbox}

%================================================
% Importance
%================================================

\booksubheading{Programming Environment কেন প্রয়োজন?}

Programming Environment-এর সাহায্যে—

\begin{compacttwocolpoints}

\item সহজে Program Code লেখা যায়।

\item Program Run বা Execute করা যায়।

\item Error শনাক্ত করা যায়।

\item Code পরিবর্তন ও সংশোধন করা যায়।

\item Program File Save করা যায়।

\item Program Development-এর কাজ সুশৃঙ্খলভাবে করা যায়।

\end{compacttwocolpoints}

%================================================
% Language vs Environment
%================================================

\booksubheading{Programming Language ও Programming Environment এক নয়}

\textbf{Programming Language} হলো Program লেখার ভাষা।

উদাহরণ—

\begin{center}

\textbf{Python}
\(\longrightarrow\)
Programming Language

\end{center}

অন্যদিকে,

\textbf{Programming Environment} হলো সেই ভাষায় Code লেখা,
Run করা ও Development-এর প্রয়োজনীয় পরিবেশ।

\begin{center}

\textbf{IDLE / VS Code / PyCharm}
\(\longrightarrow\)
Development Environment / Tool

\end{center}

%================================================
% One Look Table
%================================================

\booksubheading{এক নজরে মনে রাখবে}

\begin{tabularx}
{\linewidth}
{
|>{\bfseries\RaggedRight\arraybackslash}p{0.30\linewidth}
|Y|
}

\hline

\rowcolor{softBlue}

\multicolumn{1}{|c|}{\textbf{বিষয়}}
&
\multicolumn{1}{c|}{\textbf{কাজ}}

\\

\hline

Python
&
Programming Language
\\

\hline

Interpreter
&
Python Code Execute করে
\\

\hline

Code Editor
&
Code লেখা ও Edit করা হয়
\\

\hline

Shell
&
সরাসরি Instruction দিয়ে দ্রুত Result দেখা যায়
\\

\hline

IDE
&
Development-এর বিভিন্ন সুবিধা একসঙ্গে দেয়
\\

\hline

IDLE
&
Python শেখা, Code লেখা ও Run করার একটি পরিচিত Environment
\\

\hline

\end{tabularx}

%================================================
% Common Mistakes
%================================================

\booksubheading{সাধারণ ভুল}

\begin{compacttwocolpoints}

\item Python ও IDLE-কে একই বিষয় মনে করা।

\item IDE এবং IDLE-এর পূর্ণরূপ গুলিয়ে ফেলা।

\item Shell ও Editor-কে একই কাজের Tool মনে করা।

\item Interpreter-কে Code Editor মনে করা।

\item Programming Environment-কে Programming Language মনে করা।

\item Shell-এ কাজ করার জন্য প্রতিবার Program File তৈরি করতেই হবে
এমন ধারণা করা।

\end{compacttwocolpoints}

%================================================
% Final Reminder
%================================================

\begin{rememberbox}

\textbf{\color{deepOrange}মনে রাখবে}

\begin{itemize}

\item
Program লেখা, Edit, Run এবং Error সংশোধনের প্রয়োজনীয় পরিবেশকে
\textbf{Programming Environment} বলে।

\item
Python Program Execute করতে
\textbf{Python Interpreter}
ব্যবহৃত হয়।

\item
Code Editor-এ Program Code লেখা ও Edit করা যায়।

\item
Shell-এ Instruction দিয়ে দ্রুত Result দেখা যায়।

\item
\textbf{IDE}
=
Integrated Development Environment।

\item
\textbf{IDLE}
=
Integrated Development and Learning Environment।

\item
IDLE Python-এর পরিচিত Programming Environment।

\item
Shell ছোট Test-এর জন্য এবং Editor পূর্ণ Program লিখে
Save করার জন্য সুবিধাজনক।

\item
দ্রুত মনে রাখবে—

\begin{center}

\textbf{Python}
\(\longrightarrow\)
Language
\qquad
\textbf{Interpreter}
\(\longrightarrow\)
Execute

\medskip

\textbf{Editor}
\(\longrightarrow\)
Write/Edit
\qquad
\textbf{Shell}
\(\longrightarrow\)
Quick Result

\medskip

\textbf{IDE}
\(\longrightarrow\)
সব Development সুবিধা একসাথে

\end{center}

\end{itemize}

\end{rememberbox}

\small

%================================================
% 5.28 Python Installation and Setup
%================================================

